\section{Background}
\label{sec:background}

Most PUBG prediction studies estimate final placement from post-match
statistics~\cite{kagglePubgFinishPlacement,sen2021pubgFinalRank,lee2025pubgRank}.
That framing is appropriate for final-outcome modeling, but it removes the
temporal order and spatial context of risk. The closest prior work is Sagong
et al.'s phase-based framework, which scores player aggressiveness and
activity with FightScore and ActiveScore, clusters strategy types, and
predicts next-phase survival~\cite{sagong2025hunterCamper}. We inherit their
phase-level framing but change the unit and representation: teams rather than
players, and full telemetry reconstruction rather than a small strategy-score
summary.

The shrinking safe zone also motivates a spatio-temporal view. Habitat
fragmentation and movement ecology treat viable space, movement, exposure,
and density as coupled time-space quantities
\cite{fahrig2003fragmentation,nathan2008movement}. We use this only as an
analytic analogy: PUBG teams are not organisms, but their survival conditions
are shaped by a contracting activity space. This motivates representing each
phase as both a feature row and a graph snapshot. Graph edges encode who is
near whom, who attacks whom, and which teams share safe-zone pressure
\cite{sanchezLengeling2021gnnIntro}.
